\documentclass[12pt]{article}

\usepackage{amsmath, amsfonts}
\usepackage[margin=1cm]{geometry}

\begin{document}

\noindent
{\large\bfseries Number Theory}
\hfill
znz
\quad
\today

\hrule
\medskip

\begin{enumerate}
  \item  Suppose
        \[
        a\equiv a'\pmod n,
        \qquad
        b\equiv b'\pmod n.
        \]Prove that
        \[
        a+b\equiv a'+b'\pmod n
        \]and
        \[
        ab\equiv a'b'\pmod n.
        \]Conclude that
        \[
        [a]+[b]=[a+b],
        \qquad
        [a][b]=[ab]
        \]are well-defined operations on \(\mathbb Z/n\mathbb Z\).
  \item Consider
        \[
        \mathbb Z/8\mathbb Z.
        \]For each element
        \[
        [0],[1],\ldots,[7],
        \]determine whether it has a multiplicative inverse.
        For every invertible element, explicitly find its inverse.
        Then formulate a conjecture relating invertibility of \([a]\) to
        \[
        \gcd(a,8).
        \]Do not prove the general conjecture yet.
  \item Let \(n>1\) be composite.
        Prove directly that there exist
        \[
        [a]\neq[0],
        \qquad
        [b]\neq[0]
        \]in \(\mathbb Z/n\mathbb Z\) such that
        \[
        [a][b]=[0].
        \]Then prove that this implies \(\mathbb Z/n\mathbb Z\) cannot be a field.
        Try to write this as one clean implication chain rather than relying on an example.
  \item Let \(p\) be prime and let
        \[
        [a]\in\mathbb Z/p\mathbb Z,
        \qquad
        [a]\neq[0].
        \]Using Bézout's identity,
        \[
        ax+py=\gcd(a,p)
        \]for suitable \(x,y\in\mathbb Z\), prove that \([a]\) has a multiplicative inverse.
        Conclude:
        \[
        \boxed{
        \mathbb Z/n\mathbb Z\text{ is a field}
        \iff
        n\text{ is prime}.
        }
        \]This is today's main proof.
  \item Work in \(\mathbb F_5\).
        Compute
        \[
        1^4,\quad2^4,\quad3^4,\quad4^4
        \]modulo \(5\).
        Then formulate a conjecture for
        \[
        a^{p-1}\pmod p
        \]when \(p\) is prime and
        \[
        p\nmid a.
        \]Do not try to prove the general statement yet.
        This observation will lead to Fermat's little theorem, finite multiplicative groups, and later the Frobenius map
        \[
        x\mapsto x^p,
        \]which becomes a central actor in the function-field version of RH.
\end{enumerate}

\end{document}
